% ID: SEL-D0006-OA-E02
% ESTADO: validada. S007, 2026-11-07; tres pasadas del mismo autor y revisión matemática, didáctica y visual.
\textbf{a) Dominio y derivada.}\par\nopagebreak[4]
El denominador no puede anularse: $x+1\ne0$. Así, $\operatorname{Dom}f=\mathbb R\setminus\{-1\}$. Aplicando la regla del cociente,
\[
f'(x)=1-\frac{3(x+1)-(3x-1)}{(x+1)^2}
=\boxed{1-\frac4{(x+1)^2}},\quad x\ne-1.
\]
Comprobación alternativa: $f(x)=x-3+4/(x+1)$, cuya derivada da la misma expresión.

\textbf{b) Pendiente de la tangente.}\par\nopagebreak[4]
La pendiente en una abscisa es el valor de la derivada allí; $2/3$ pertenece al dominio:
\[
m=f'(2/3)=1-\frac4{(5/3)^2}=1-\frac{36}{25}=\boxed{-\frac{11}{25}=-0.44}.
\]

\textbf{c) Tangentes horizontales.}\par\nopagebreak[4]
Una tangente horizontal exige $f'(x)=0$ y un punto del dominio:
\[
1-\frac4{(x+1)^2}=0\iff (x+1)^2=4\iff x=1\text{ o }x=-3.
\]
Ambas abscisas son admisibles. Calculamos las ordenadas: $f(1)=1-2/2=0$ y $f(-3)=-3-(-10)/(-2)=-8$. Los puntos son
\[
\boxed{(1,0)\quad\text{y}\quad(-3,-8).}
\]
Las rectas tangentes respectivas son $y=0$ e $y=-8$. Sustituyendo ambas abscisas en la derivada se obtiene exactamente cero; no se incluye $x=-1$, donde no existe la función.
